% Part: lambda-calculus
% Chapter: church-rosser
% Section: beta-eta-reduction

\documentclass[../../../include/open-logic-section]{subfiles}

\begin{document}

\olfileid[id]{lam}{cr}{be}

\olsection{Reduksi-$\beta\eta$}

Sifat Church--Rosser berlaku bagi reduksi-$\beta\eta$ ($\bered$).

\begin{lem}\ollabel{lem:one-par}
  Jika $M \beredone M'$, maka $M \beredpar M'$.
\end{lem}

\begin{proof}
  Dengan induksi pada !!{derivation} $M \beredone M'$. Jika $M \beredone M'$
  melalui konversi-$\eta$ (yakni \olref[syn][eta]{defn:beredone}), kita
  menggunakan \olref[pbe]{thm:refl}. Kasus lainnya sama seperti dalam
  \olref[b]{lem:one-par}.
\end{proof}

\begin{lem}\ollabel{lem:par-red}
  Jika $M \beredpar M'$, maka $M \bered M'$.
\end{lem}

\begin{proof} Dengan induksi pada !!{derivation} $M \beredpar M'$.

  Jika aturan terakhir adalah \olref[pbe]{defn:beredpar5}, maka
  $M=\lambd[x][Nx]$ dan $M'=N'$ untuk suatu $x$, $N$, dan~$N'$ dengan
  $x\notin\FV{N}$ serta $N\beredpar N'$. Kita dapat terlebih dahulu mereduksi
  $\lambd[x][Nx]$ menjadi $N$ melalui konversi-$\eta$, lalu mengikuti
  rangkaian langkah $\beredone$ yang menunjukkan $N\bered N'$, yang berlaku
  berdasarkan hipotesis induksi.
\end{proof}

\begin{lem}\ollabel{lem:str}
  $\bered$ adalah relasi transitif terkecil yang memuat $\beredpar$.
\end{lem}

\begin{proof}
  Sama seperti dalam \olref[b]{lem:str}.
\end{proof}

\begin{thm}\ollabel{thm:cr}
  $\bered$ memenuhi sifat Church--Rosser.
\end{thm}

\begin{proof}
  Berdasarkan \olref[dap]{thm:str}, \olref[pbe]{thm:cr}, dan
  \olref{lem:str}.
\end{proof}
\end{document}
